% ============================================================================
% Fichier  : partie0/P00_certifiee_compromise.tex
% Titre    : Certifiée, et pourtant Compromise
% Auteur   : MINKA MI NGUIDJOÏ Thierry Emmanuel
% Version  : 1.0 -- novembre 2026
% Corrigés : corriges/Corr_P00.tex
% ============================================================================
% Faits du cas (bible des cas, Claude outputs/02_bible_des_cas.md) :
% - corpus : BANQUE SANAGA, Douala, 38 agences, DG M. FOTSO, COBAC ; bascule
%   ratée du système central en mars 2026 (GCI P00) ;
% - propres à SEC08 : certificat ISO/IEC 27001:2022 de juin 2025 (CERTIQUAL
%   AFRIQUE, organisme fictif) sur le seul centre de données principal de
%   Douala ; rançongiciel du 14/09/2026 entré par l'agence de Ngaoundéré,
%   180 postes dans 9 agences, 3 jours, rançon de 90 M FCFA non payée, coût
%   410 M FCFA ; Mme ABENA, RSSI.
% ISO/IEC 27001:2022, article 4.3 : paraphrasé (texte sous licence), jamais
% cité.
% ============================================================================

\chapter[Certifiée, et pourtant Compromise]{Prologue : Certifiée, et pourtant Compromise}
\label{chap:prologue}
\chaptermark{Prologue}

\epigraph{%
    « Un certificat dit ce qui a été audité.\\
    Il ne dit rien de ce qui ne l'a pas été. »%
}{Principe du cours}

\niveauChapitre{Introductif}{%
    Aucun prérequis de gestion des risques ni d'audit. Les notions de base de
    la cybersécurité (MSN-1S1-SEC01) suffisent.%
}

\begin{boxobjectifs}
\begin{itemize}
    \item Reconstituer, à partir de ses traces, la compromission d'une banque
          certifiée ISO/IEC~27001.
    \item Distinguer un certificat, la conformité à une norme et la sécurité
          effective.
    \item Lire le périmètre d'un système de management de la sécurité et
          repérer un périmètre de complaisance.
    \item Reconnaître qu'un risque jamais apprécié n'est pas un risque faible,
          mais une inconnue.
    \item Lire un ensemble de contrôles par trois mesures~: $\mu$, $U$ et $H$.
    \item Énoncer les cinq réflexes du cours et situer ses parties, ses
          ateliers et ses cas.
\end{itemize}
\end{boxobjectifs}

\bigskip

% ============================================================================
\section*{Avant de commencer : un certificat au mur, des guichets fermés}
% ============================================================================

Douala, jeudi 17~septembre 2026. Dans le hall du siège de la BANQUE
SANAGA\footnote{La BANQUE SANAGA, ses dirigeants, ses prestataires,
l'organisme certificateur, les dates et les montants de ce cas sont fictifs.
Les mécanismes d'attaque, les normes et les textes cités sont exacts.}, un
cadre doré présente le certificat ISO/IEC~27001 obtenu quinze mois plus tôt.
Dans les étages, la cellule de crise en est à sa troisième nuit. Neuf
agences sur trente-huit ont gardé leurs guichets fermés pendant trois jours~;
180 postes de travail ont été chiffrés par un rançongiciel~; les attaquants
réclament 90~millions de francs CFA.

Le directeur général, M.~FOTSO, pose la question que tout le monde se pose~:
« Nous sommes certifiés. Comment est-ce possible~? » La responsable de la
sécurité des systèmes d'information, Mme~ABENA, répond par une seule
phrase~: « Le certificat porte sur le centre de données. L'attaque est
passée par une agence. »

\medskip

Ce prologue reconstitue l'incident, montre pourquoi il ne contredit pas le
certificat, et en tire les réflexes qui structurent tout le cours. Il ne
suppose aucune connaissance préalable des normes ni de l'audit~: les notions
sont nommées ici et définies dans les chapitres qui suivent.

\begin{boxalert}
\textbf{Violation typique.} Présenter un certificat ISO/IEC~27001 comme la
preuve que « la banque est sécurisée ».

\smallskip
\textbf{Règle~:} un certificat atteste qu'un système de management, sur un
périmètre donné, satisfait les exigences d'une norme à la date de l'audit.
On lit toujours le périmètre avant de lire le certificat.
\end{boxalert}

\separateur

% ============================================================================
\section{Ce qui s'est passé}
\label{sec:pro-faits}
% ============================================================================

\subsection{La chronologie}

La RSSI a reconstitué l'incident à partir des journaux des postes, de la
messagerie et du partage de fichiers, puis l'a présenté au comité des
risques. Chaque ligne du Tableau~P.1 renvoie à une trace.

\begin{tableaucours}{L{2.8cm}L{6.4cm}L{3.0cm}}{%
    \tableauHeader{Date} &
    \tableauHeader{Événement} &
    \tableauHeader{Trace}
}
Juin 2025 & Certificat ISO/IEC~27001:2022 délivré par CERTIQUAL AFRIQUE pour
trois ans. Périmètre~: le centre de données principal de Douala &
Certificat et rapport d'audit de certification \\
Novembre 2025 & Revue annuelle des risques~: aucun scénario ne porte sur le
réseau des agences, hors périmètre & Registre des risques 2025 \\
Lundi 14 sept. 2026, 10~h~12 & Un agent de l'agence de Ngaoundéré ouvre la
pièce jointe d'un faux avis de virement & Journal de la messagerie \\
14 sept., 10~h~14 & Un programme malveillant s'installe sur le poste~; il
récupère les identifiants enregistrés dans le navigateur & Journal du
poste \\
14 -- 15 sept. & Avec ces identifiants, l'attaquant accède au partage de
fichiers commun aux agences et s'y déplace & Journal du serveur de
fichiers \\
Mardi 15 sept., 2~h~40 & Chiffrement simultané de 180 postes dans 9
agences~; message de rançon & Journaux des postes \\
15 -- 17 sept. & Guichets fermés dans les 9 agences~; reconstruction des
postes à partir d'images propres & Rapport de la cellule de crise \\
18 sept. & Réouverture des agences~; la rançon n'est pas payée & Décision
du directeur général \\
\end{tableaucours}
\vspace{-0.8em}
\begin{center}{\small\itshape Tableau~P.1. Chronologie de l'incident
(reconstitution de la RSSI)}\end{center}

Le centre de données certifié n'a pas été atteint~: il n'était pas relié au
partage de fichiers des agences. Les comptes, les opérations et le système
central sont restés intacts. Le coût de l'incident, chiffré par la
direction financière, atteint pourtant \textbf{410~millions de francs CFA}~:
arrêt des guichets, reconstruction des postes, prestataire de réponse à
incident, communication auprès de la clientèle.

\subsection{Constater, puis analyser}

Comme dans toute investigation, on sépare ce qui est établi de ce que l'on
en déduit.

\begin{boxscenario}{Opération SANAGA. Le rançongiciel}
\textbf{Constatations.}
\begin{itemize}
    \item Le certificat couvre le seul centre de données principal de
          Douala~; les 38 agences, la messagerie et les postes de travail sont
          hors périmètre.
    \item Le registre des risques 2025 ne contient aucun scénario sur les
          agences.
    \item L'attaque est entrée par un courriel ouvert en agence, puis s'est
          propagée par un partage de fichiers commun aux agences.
    \item Les identifiants volés n'étaient protégés par aucun second facteur
          d'authentification.
\end{itemize}

\textbf{Analyse} (confiance élevée). Le certificat n'est pas faux~: le
système de management certifié couvre bien le centre de données, et celui-ci
n'a pas été touché. Ce qui a échoué, c'est le choix du périmètre et
l'appréciation des risques, qui s'arrêtait à ses frontières. La banque a
protégé ce qu'elle avait décidé d'auditer, et laissé dans l'ombre ce qu'elle
n'avait pas regardé.
\end{boxscenario}

% ============================================================================
\section{Certificat, conformité, sécurité}
\label{sec:pro-certificat}
% ============================================================================

\subsection{Trois choses différentes}

\begin{boxdef}
\termecle{Système de management de la sécurité de
l'information}\index{SMSI} (\sigle{SMSI})~: ensemble de politiques, de
processus, de rôles et de contrôles par lequel une organisation pilote la
sécurité de son information sur un périmètre défini, selon un cycle
d'amélioration continue. L'ISO/IEC~27001 en fixe les
exigences~\cite{iso27001_2022}.

\medskip
\termecle{Certificat}\index{certificat}~: attestation délivrée par un
organisme tiers accrédité, après un audit, qu'un système de management
satisfait les exigences d'une norme, sur un périmètre et pour une durée
donnés.

\medskip
\termecle{Conformité}\index{conformité}~: satisfaction d'une exigence
précise (d'une norme, d'une loi, d'un contrat), établie par une preuve.

\medskip
\termecle{Sécurité effective}\index{sécurité effective}~: capacité réelle de
l'organisation à résister aux attaques qu'elle subit, et à s'en relever.
\end{boxdef}

\begin{tableaucours}{L{2.6cm}L{4.4cm}L{5.2cm}}{%
    \tableauHeader{Notion} &
    \tableauHeader{Ce qu'elle établit} &
    \tableauHeader{Ce qu'elle n'établit pas}
}
Certificat & Un SMSI conforme à la norme sur le périmètre, à la date de
l'audit & La sécurité hors périmètre~; la sécurité après l'audit \\
Conformité & Qu'une exigence précise est satisfaite & Que l'exigence suffit à
couvrir le risque \\
Sécurité effective & Que les attaques réelles échouent ou sont contenues &
Rien sans mesure~: elle se démontre par des tests et des incidents \\
\end{tableaucours}
\vspace{-0.8em}
\begin{center}{\small\itshape Tableau~P.2. Trois notions à ne pas
confondre}\end{center}

\subsection{Le périmètre, première décision de sécurité}

L'ISO/IEC~27001 demande à l'organisation de fixer les limites de son SMSI en
tenant compte de son contexte, des attentes des parties intéressées, et des
interfaces et dépendances entre ses propres activités et celles des
autres~\cite{iso27001_2022}. Le périmètre n'est donc pas libre~: il doit être
défendable. Celui de la BANQUE SANAGA avait été choisi pour obtenir
rapidement le certificat, en excluant ce qui était difficile à maîtriser.

\begin{boxdef}
\termecle{Périmètre de complaisance}\index{périmètre de complaisance}~:
périmètre de SMSI choisi pour faciliter la certification plutôt que pour
couvrir les risques de l'organisation, au point d'exclure des activités ou
des actifs dont dépend le périmètre retenu.
\end{boxdef}

Les agences de la banque n'étaient pas un sujet à part~: elles reçoivent les
clients, saisissent les opérations et sont reliées au système central. Un
attaquant installé en agence se trouve à quelques pas de ce que le
certificat protège. Le Chapitre~\ref{chap:smsi} montrera comment construire
un périmètre défendable, et comment l'auditeur le conteste.

\begin{boxinfo}
\textbf{Et le régulateur~?} Le règlement COBAC R-2016/04 impose aux
établissements de crédit un dispositif de contrôle interne qui couvre
l'ensemble de leurs risques, y compris les risques opérationnels liés aux
systèmes d'information~\cite{cobac2016r04}. Un certificat sur une partie du
système ne dispense pas la banque de maîtriser l'autre. Le
Chapitre~\ref{chap:droit} présente ces obligations.
\end{boxinfo}

% ============================================================================
\section{Un Risque jamais Apprécié n'est pas un Risque Faible}
\label{sec:pro-risque}
% ============================================================================

\subsection{Le trou dans le registre}

Le registre des risques de 2025 recensait des scénarios sur le centre de
données~: panne électrique, défaillance du prestataire d'hébergement,
intrusion par Internet. Aucun ne portait sur les agences. Lors de la revue,
le comité des risques a lu le registre et conclu qu'aucun risque
intolérable n'était ouvert.

Le raisonnement est faux, et c'est l'erreur la plus répandue en gestion des
risques~: un risque absent du registre a été lu comme un risque nul. Il
n'était pas nul~; il était \emph{inconnu}. La bonne cote n'était pas
« faible », mais « non évalué », et cette cote aurait dû déclencher une
appréciation.

\begin{boxalert}
\textbf{Violation typique.} Lire l'absence d'un scénario dans le registre
comme la preuve que le risque est maîtrisé.

\smallskip
\textbf{Règle~:} ce qui n'a pas été apprécié est coté \NEK{} (connaissance
insuffisante), pas « faible ». Une cote de vraisemblance sans preuve est une
opinion~; elle se signale comme telle.
\end{boxalert}

\subsection{Une fiche de risque}

Voici ce qu'aurait dû contenir le registre. Le Chapitre~\ref{chap:risque}
en définira chaque rubrique, et le Chapitre~\ref{chap:ebios} montrera comment
construire le scénario avec la méthode EBIOS Risk Manager.

\begin{boxobjet}{R-AG-01 · Rançongiciel en agence}
\textbf{Actif essentiel~:} le service au guichet des agences (versements,
retraits, ouvertures de compte).

\textbf{Actifs supports~:} postes de travail des agences, messagerie,
partage de fichiers commun, comptes des agents.

\textbf{Source de risque~:} groupe cybercriminel motivé par le gain.

\textbf{Événement redouté~:} indisponibilité des guichets pendant plusieurs
jours dans plusieurs agences.

\textbf{Scénario~:} courriel piégé ouvert en agence, vol d'identifiants,
déplacement par le partage commun, chiffrement des postes.

\textbf{Vraisemblance~:} \NEK{} en 2025 (scénario jamais apprécié)~; après
l'incident, élevée.

\textbf{Impact~:} perte d'exploitation, coûts de reconstruction, atteinte à la
réputation, compte rendu au régulateur.

\textbf{Mesures existantes~:} antivirus sur les postes~; aucun second facteur~;
partage de fichiers ouvert à toutes les agences~; sauvegardes des postes non
testées.

\textbf{Risque résiduel~:} intolérable, tant que le partage n'est pas
cloisonné et que l'authentification n'est pas renforcée.
\end{boxobjet}

% ============================================================================
\section{Lire la Sécurité par Trois Mesures}
\label{sec:pro-mesures}
% ============================================================================

\subsection{Deux colonnes, un certificat}

Après l'incident, la direction de l'audit interne, dirigée par
Mme~EKAMBI, évalue dix contrôles clés de sécurité de deux façons~: sur le
centre de données certifié, puis sur le réseau des agences. Chaque contrôle
est coté \OUI{} (en place et prouvé), \NON{} (absent ou défaillant) ou
\NEK{} (impossible à établir faute de preuve).

\begin{tableaucours}{L{6.6cm}cc}{%
    \tableauHeader{Contrôle} &
    \tableauHeader{Centre de données} &
    \tableauHeader{Agences}
}
Second facteur d'authentification pour les accès à distance & \OUI{} & \NON{} \\
Correctifs de sécurité appliqués dans le délai fixé & \OUI{} & \NEK{} \\
Sauvegardes hors ligne testées & \OUI{} & \NON{} \\
Cloisonnement du réseau & \OUI{} & \NON{} \\
Comptes d'administration nominatifs et revus & \OUI{} & \NEK{} \\
Journalisation centralisée et revue & \OUI{} & \NEK{} \\
Filtrage des pièces jointes de la messagerie & \OUI{} & \OUI{} \\
Sensibilisation des agents dans l'année & \NEK{} & \NEK{} \\
Plan de réponse aux incidents testé & \NON{} & \NON{} \\
Protection des postes (antivirus) & \OUI{} & \OUI{} \\
\textbf{Total (\OUI{} / \NON{} / \NEK{})} & \textbf{8 / 1 / 1} & \textbf{2 / 4 / 4} \\
\end{tableaucours}
\vspace{-0.8em}
\begin{center}{\small\itshape Tableau~P.3. Dix contrôles clés, deux
périmètres}\end{center}

\subsection{Trois mesures}

Le cours lit un ensemble de cotes par trois mesures, définies au
Chapitre~\ref{chap:conclure}~:
\begin{itemize}
    \item le \termecle{taux de conformité}\index{taux de conformité}~$\mu$, la
          part des \OUI{}~;
    \item le \termecle{taux d'incertitude}\index{taux d'incertitude}~$U$, la
          part des \NEK{}~;
    \item l'\termecle{entropie}\index{entropie}~$H$, la dispersion des
          réponses, en bits, de 0 à $\log_2 3 \approx 1{,}585$~:
          $H = -\sum_m p(m)\log_2 p(m)$.
\end{itemize}

\begin{boxlecture}{Deux périmètres de la BANQUE SANAGA}
\textbf{Constatation.}
\begin{itemize}
    \item Centre de données (8/1/1)~: $\mu = 0{,}80$, $U = 0{,}10$,
          $H = -(0{,}8\log_2 0{,}8 + 2 \times 0{,}1\log_2 0{,}1) \approx 0{,}92$~bit.
    \item Agences (2/4/4)~: $\mu = 0{,}20$, $U = 0{,}40$,
          $H = -(0{,}2\log_2 0{,}2 + 2 \times 0{,}4\log_2 0{,}4) \approx 1{,}52$~bit.
    \item Ensemble (10/5/5)~: $\mu = 0{,}50$, $U = 0{,}25$, $H = 1{,}50$~bit.
\end{itemize}

\textbf{Analyse} (confiance élevée sur le calcul, moyenne sur la lecture).
Le certificat décrit la première ligne seulement~: un périmètre largement
conforme et bien connu. Les agences combinent une conformité faible et une
forte incertitude~: quatre contrôles sur dix ne peuvent pas même être
évalués, faute de trace. C'est l'endroit où l'attaque est entrée. Lue
d'un bloc, la banque paraît « à moitié conforme »~; lue par périmètre, elle
montre un îlot maîtrisé dans une zone inconnue.

\medskip
Une cote mérite enfin d'être relue~: le filtrage des pièces jointes est
coté \OUI{} en agence, alors que l'attaque est entrée par une pièce jointe.
Il n'y a pas de contradiction~: \OUI{} dit que le contrôle est en place et
prouvé, pas qu'il est infaillible. Un filtre arrête les pièces jointes
connues comme malveillantes~; celle du 14~septembre ne l'était pas encore.
C'est le troisième réflexe du cours~: une mesure se juge à son
fonctionnement face aux attaques réelles.

\medskip
\hfill\lecture{0{,}20}{0{,}40}{1{,}52}
\end{boxlecture}

\begin{boxinfo}
\textbf{Une règle à retenir dès maintenant.} On ne conclut jamais sur une
seule mesure. Un taux de conformité de 0,50 ne dit pas si la moitié du
système est solide et l'autre fragile, ou si l'on ignore l'état de la
moitié du système~: $U$ et $H$ le disent. Le Chapitre~\ref{chap:conclure}
en tirera la règle de décision~: $U$ d'abord, $\mu$ ensuite, $H$ pour
contrôler.
\end{boxinfo}

% ============================================================================
\section{Cinq Réflexes}
\label{sec:pro-reflexes}
% ============================================================================

L'incident se résume en cinq réflexes. Chacun fait l'objet d'au moins un
chapitre.

\begin{tableaucours}{cL{4.0cm}L{6.0cm}c}{%
    \tableauHeader{N°} &
    \tableauHeader{Réflexe} &
    \tableauHeader{Ce qu'il évite} &
    \tableauHeader{Chapitre}
}
1 & Lire le périmètre avant le certificat & Croire protégé ce qui n'a jamais
été audité & \ref{chap:smsi} \\
2 & Ce qui n'est pas apprécié est inconnu, pas faible & Un registre qui
rassure parce qu'il est incomplet & \ref{chap:risque} \\
3 & Une mesure se prouve en fonctionnement & Des contrôles qui existent sur
le papier & \ref{chap:audit} \\
4 & Mesurer l'incertitude & Une conformité moyenne qui cache des zones
inconnues & \ref{chap:conclure} \\
5 & Relier norme, risque et audit & Trois exercices séparés qui ne se
parlent pas & \ref{chap:grc} \\
\end{tableaucours}
\vspace{-0.8em}
\begin{center}{\small\itshape Tableau~P.4. Les cinq réflexes du
cours}\end{center}

\begin{boxprocedure}
\textbf{Le minimum d'une cote documentée.} Pour qu'une cote de contrôle ou
de risque puisse être défendue devant un comité, un auditeur ou le
régulateur, le papier de travail qui la porte indique au moins~:
\begin{itemize}
    \item le contrôle ou le scénario concerné, et son périmètre~;
    \item le critère (la norme, la politique ou la loi qui fixe l'attendu)~;
    \item la cote (\OUI{}, \NON{} ou \NEK{}) et la preuve qui l'appuie~;
    \item pour un \NEK{}, la raison pour laquelle on ne peut pas savoir~;
    \item la date, l'auteur et le réviseur.
\end{itemize}
\end{boxprocedure}

\subsection{Ce que l'incident dit à chaque métier}

\begin{boxmetier}{RSSI}
\textbf{Ce que Mme~ABENA retient}

Le périmètre du SMSI était une décision de sécurité, et non une formalité
de certification. Elle propose au comité d'étendre le SMSI aux agences en
deux ans, de cloisonner le partage de fichiers et d'imposer un second
facteur pour tous les accès. Elle sait que le prochain audit de
surveillance de l'organisme certificateur ne lui sera d'aucun secours sur
ces sujets, puisqu'ils sont hors périmètre.
\end{boxmetier}

\begin{boxmetier}{Auditeur en cybersécurité}
\textbf{Contester le périmètre}

L'auditeur ne se contente pas de vérifier ce qui est inclus~: il examine
ce qui est exclu, et la justification de l'exclusion. Un périmètre qui
exclut une activité dont dépend le périmètre retenu est une constatation
d'audit en soi (Chapitre~\ref{chap:smsi}).
\end{boxmetier}

\begin{boxmetier}{Gestionnaire des risques cyber}
\textbf{Apprécier hors des frontières}

L'appréciation des risques ne s'arrête pas au périmètre du certificat~: elle
porte sur les activités de l'organisation. Le gestionnaire des risques ajoute
le scénario R-AG-01, le chiffre (Chapitre~\ref{chap:fair}) et le présente
au comité avec ses options de traitement.
\end{boxmetier}

\begin{boxmetier}{Responsable GRC}
\textbf{Un seul dispositif}

Le registre des risques, la déclaration d'applicabilité du SMSI et le
programme d'audit doivent couvrir le même univers. Le responsable GRC
vérifie qu'aucun pan de l'activité n'échappe aux trois à la fois
(Chapitre~\ref{chap:grc-pratique}).
\end{boxmetier}

\begin{boxmetier}{Consultant en mise en œuvre de SMSI}
\textbf{Conseiller un périmètre défendable}

Le consultant qui a accompagné la certification a proposé le périmètre le
plus facile à certifier. Son rôle était aussi de dire ce que ce choix
laissait sans protection. Un périmètre étroit peut être une étape~; il doit
alors être présenté comme tel, avec un calendrier d'extension.
\end{boxmetier}

% ============================================================================
\section{Le Cours et son Fil Rouge}
\label{sec:pro-cours}
% ============================================================================

\subsection{Une discipline, une mesure}

Ce cours enseigne à appliquer les normes de sécurité, à gérer les risques
et à auditer la cybersécurité. Il les relie par un même fil rouge~: la
mesure de ce que l'on sait et de ce que l'on ignore. L'instrument de
réponse \OUI{}~/ \NON{}~/ \NEK{} et les mesures $\mu$, $U$ et $H$ viennent du
programme de recherche de l'auditologie~\cite{minka2026auditologie,minka2026nek}.
Le cours les applique à l'audit, comme les autres cours du corpus, et aussi
à l'appréciation des risques, où une vraisemblance cotée sans preuve est un
\NEK{} déguisé.

\begin{boxinfo}
\textbf{Ce que dit la recherche.} L'article \emph{Vers l'auditologie} pose
l'\emph{observabilité partielle} comme premier axiome~: l'auditeur ne voit
jamais tout le système, il en voit des traces. Le périmètre d'un certificat
en est une illustration~: il délimite ce qui a été observé. L'article sur
l'incertitude épistémique en tire l'instrument trichotomique et ses
mesures. Les travaux sur la décomposition, le suivi dans le temps et le jeu
entre auditeur et audité sont en préparation~; le Chapitre~\ref{chap:auditologie}
les présente comme tels.
\end{boxinfo}

\subsection{Trois cas}

\begin{tableaucours}{L{3.2cm}L{5.2cm}L{3.8cm}}{%
    \tableauHeader{Cas} &
    \tableauHeader{Situation} &
    \tableauHeader{Thèmes dans ce cours}
}
BANQUE SANAGA (Douala) & Banque de 38 agences, certifiée ISO/IEC~27001 sur
son centre de données, compromise en agence & SMSI, périmètre,
quantification, régulation bancaire \\
LiZbiZ (Douala) & Logistique, quatre centres de données, douze agences,
écart de stock et fraude & EBIOS Risk Manager, audit technique,
décomposition \\
ANPP (Yaoundé) & Agence publique de paiement & Audit de sécurité
réglementaire, correspondances de référentiels \\
\end{tableaucours}
\vspace{-0.8em}
\begin{center}{\small\itshape Tableau~P.5. Les trois cas fil rouge}\end{center}

\subsection{Le parcours}

\begin{tableaucours}{L{3.8cm}cL{6.8cm}}{%
    \tableauHeader{Partie} &
    \tableauHeader{Chapitres} &
    \tableauHeader{Question traitée}
}
I. Fondements & 1 à 2 & Qu'est-ce que gouverner la sécurité~? Qu'est-ce
qu'un risque cyber, et que sait-on de lui~? \\
II. Méthodes & 3 à 7 & Comment apprécier, traiter et chiffrer un risque~?
Comment conduire un audit et conclure~? \\
III. Cadre normatif & 8 à 10 & Quelles normes, quels référentiels, quel
droit, quelle déontologie~? \\
IV. Pratique & 11 à 14 & Comment mettre en œuvre et auditer un SMSI,
auditer la technique, faire vivre la GRC~? \\
V. Défis avancés & 15 & Tiers, nuage, intelligence artificielle, audit
continu \\
VI. Recherche & 16 & Mesurer ce que l'on ne sait pas \\
\end{tableaucours}
\vspace{-0.8em}
\begin{center}{\small\itshape Tableau~P.6. Les six parties du cours}\end{center}

Le cours compte 30~heures en présentiel, réparties en quinze séances de
deux heures qui couvrent le prologue et les six parties. Les fiches de
questions-réponses de chaque chapitre et les dix ateliers d'environ
90~minutes de NDAP SUKULU, sous le code \codeUE{} (registre des risques,
EBIOS Risk Manager en deux séances, quantification FAIR, décomposition et
cotation, déclaration d'applicabilité, table de correspondance des
référentiels, audit d'un SMSI, audit technique, synthèse), sont des
travaux à faire à la maison, hors de ces 30~heures. Les corrigés des
exercices de fin de chapitre sont réunis dans un livret séparé.

\begin{boxinfo}
\textbf{Ce que le cours ne fait pas.} Il ne refait pas les autres
enseignements du master~: le test d'intrusion relève de l'UE
MSN-2S3-SEC03, la forensique de l'UE MSN-2S3-SEC07, le droit du numérique et
des données personnelles de l'UE MSN-2S3-JUR01, l'audit de protection des
données de l'UE MSN-2S4-SEC02. Il y renvoie chaque fois que nécessaire.
\end{boxinfo}

\separateur

\begin{boxresume}
\textbf{Ce qu'il faut retenir de ce prologue~:}
\begin{itemize}
    \item \textbf{Un certificat dit ce qui a été audité}~: un SMSI conforme,
          sur un périmètre, à une date. Il ne dit rien du reste.
    \item \textbf{Certificat, conformité, sécurité effective} sont trois
          notions distinctes~; seule la dernière se mesure aux attaques.
    \item \textbf{Le périmètre est une décision de sécurité}~: un périmètre
          de complaisance exclut ce dont dépend ce qu'il protège.
    \item \textbf{Un risque jamais apprécié est inconnu, pas faible}~: il se
          cote \NEK{} et appelle une appréciation.
    \item \textbf{Trois mesures}, $\mu$, $U$ et $H$, lisent un ensemble de
          contrôles~; lues par périmètre, elles montrent où se cache
          l'incertitude.
    \item \textbf{Cinq réflexes}~: lire le périmètre, coter l'inconnu, prouver
          le fonctionnement, mesurer l'incertitude, relier norme, risque et
          audit.
\end{itemize}
\end{boxresume}

\begin{boxinfo}
\textbf{Points de contrôle.} Vous avez compris ce prologue si vous savez~:
\begin{itemize}
    \item expliquer pourquoi l'incident ne contredit pas le certificat~;
    \item dire ce qu'est un périmètre de complaisance, sur un exemple~;
    \item justifier la cote \NEK{} d'un scénario jamais apprécié~;
    \item calculer $\mu$, $U$ et $H$ pour 2, 4 et 4 réponses.
\end{itemize}
\end{boxinfo}

% ============================================================================
\section*{Exercices}
% ============================================================================

\begin{boxexo}[title={\ding{45}\quad Exercice~P.1. Certificat, conformité ou sécurité~? \niveaubase}]
Pour chacune des phrases suivantes, dites si elle porte sur le certificat,
la conformité ou la sécurité effective, et si elle est exacte.
\begin{enumerate}[(a)]
    \item « Notre certificat ISO/IEC~27001 prouve que nos agences sont
          protégées. »
    \item « L'audit de certification a constaté que la politique de sauvegarde
          du centre de données était appliquée. »
    \item « Aucun des trois exercices de simulation d'attaque de l'année n'a
          permis d'atteindre le système central. »
    \item « La politique de mots de passe satisfait l'exigence de longueur de
          la politique de sécurité du groupe. »
\end{enumerate}
\end{boxexo}

\begin{boxexo}[title={\ding{45}\quad Exercice~P.2. Lire un périmètre \niveaumoyen}]
Le périmètre d'un SMSI est ainsi rédigé~: « La gestion du système
d'information du centre de données principal de Douala, à l'exclusion des
postes de travail et de la messagerie. »
\begin{enumerate}[(a)]
    \item Citez trois activités de la banque dont dépend ce périmètre et qui
          en sont exclues.
    \item Pour chacune, décrivez un scénario d'attaque qui, partant de
          l'activité exclue, atteindrait le périmètre.
    \item Proposez une rédaction plus défendable, et un calendrier
          d'extension réaliste.
\end{enumerate}
\end{boxexo}

\begin{boxexo}[title={\ding{45}\quad Exercice~P.3. Mesurer deux périmètres \niveaumoyen}]
Six mois après l'incident, la même évaluation des dix contrôles donne, pour
les agences, 6~\OUI{}, 3~\NON{} et 1~\NEK{}.
\begin{enumerate}[(a)]
    \item Calculez $\mu$, $U$ et $H$ pour les agences.
    \item Comparez avec la situation de septembre (2/4/4). Qu'est-ce qui a
          le plus changé~: la conformité ou la connaissance~?
    \item Le nombre de \NON{} est passé de 4 à 3 seulement. Est-ce une
          mauvaise nouvelle~? Justifiez.
\end{enumerate}
\end{boxexo}

\begin{boxexo}[title={\ding{45}\quad Exercice~P.4. La note au conseil d'administration \niveauexpert}]
Le conseil d'administration demande à Mme~ABENA une note d'une page.
\begin{enumerate}[(a)]
    \item Expliquez, sans jargon, pourquoi le certificat n'a pas empêché
          l'incident, et pourquoi il reste valable.
    \item Rédigez une fiche de risque complète pour le scénario « vol
          d'identifiants de messagerie d'un dirigeant », sur le modèle de
          R-AG-01.
    \item Proposez trois décisions au conseil, avec leur coût approximatif
          en ordre de grandeur et le risque qu'elles traitent.
    \item Indiquez ce que le prochain audit devra vérifier pour que le
          conseil sache si ces décisions ont été appliquées.
\end{enumerate}
\end{boxexo}

\bigskip

\begin{boxinfo}
\textbf{Pour aller plus loin.}
\begin{itemize}
    \item L'ISO/IEC~27001:2022~\cite{iso27001_2022}, articles~4.1 à 4.3~:
          le contexte de l'organisation et le périmètre du SMSI.
    \item L'article \emph{Vers l'auditologie}~\cite{minka2026auditologie},
          sur l'observabilité partielle.
    \item L'article sur l'incertitude épistémique~\cite{minka2026nek}, qui
          définit la réponse \NEK{} et les mesures $\mu$, $U$, $H$.
\end{itemize}
\end{boxinfo}

\begin{boxinfo}
\textbf{Transition.} Le prologue a montré une banque qui gérait la sécurité
de son centre de données et ignorait celle de ses agences. Avant d'apprendre
les méthodes, il faut définir ce que signifie gouverner la cybersécurité, et
qui en porte la responsabilité. C'est l'objet du Chapitre~\ref{chap:grc}.
\end{boxinfo}
